%!TEX root = ./../../thesis-trier.tex

Bondar et al.~\cite{BondarCV23} introduced
a graph construction to characterize complete reachability
that we shortly present here (but in a condensed and slightly modified form).
Another slightly different construction appeared in~\cite{BondarV18}
for the same purpose.
For a more detailed discussion we refer to~\cite{BondarCV23}.

%
%
% E_1 singleton set
% induktiv E_k vereinigung über die in SCC, das als knoten; immer mengen als knoten
%
% bild dazu mit minipage
%
% Q_2 = \{ vereingung über SCC \}  { \bigcup_{c \in S} c \mid S ist SCC }
%
% 
% E_0 nur singeltons

\IN{Bondar-Volkov-Casas graphs}
\begin{definition}[Bondar et al.~\cite{BondarV18,BondarCV23}] \label{sync_max:def:BVCgraphs}  \label{def:Gammakgraphs}
	%In~\cite{} a sequence of directed
	%graphs was associated with $\mathcal A$ that yields a sufficient and necessary criterion
	%for complete reachablity. 
	Let $\semiautomat{A}{Q}{\Sigma}{\delta}$ and set $n = |Q|$. Define the digraph $\Gamma_0(\mathcal A)$ 
	with vertex set $Q_0 = \{ \{q\} \mid q \in Q \}$ and no edges. Then, for $0 < k < n$
	define, inductively, digraphs $\Gamma_k(\mathcal A) = (Q_k, E_k)$ with
	vertices $Q_k \subseteq \mathcal P(Q)$ and edges $E_{\le k} \subseteq Q_k \times Q_k$
	by setting\footnote{The union over all sets inside a set~$S$ is sometimes also written 
as $\bigcup S$ instead of $\bigcup_{A \in S} A$, especially in axiomatic treatments of set theory w.r.t.~the axiom of union
from Zermelo–Fraenkel set theory.}
	%\begin{align*}
	%	Q_k = & \left\{ \bigcup_{A \in S} A \mid S \text{ is a strongly connected component in } E_{k-1} \right\}, \\ 
	%	E_{ \le k} = & \{ (C, C') \in Q_k \times Q_k \mid C \ne C',~  C \supseteq \excl{w}, 
	%	\\ & \qquad
	%	C' \cap \dupl{w} \ne \emptyset \text{ for some } w \in \Sigma^* \text{ with } \defect{w} = k \} 
	%	\\ & \qquad \qquad \cup \{ (C, C') \in Q_k \times Q_k \mid 
	%	\\ & \qquad \qquad \qquad \exists (D, D') \in E_{k-1} : D \cap C \ne \emptyset \land 
	%	D' \cap C' \ne \emptyset \}.
	%\end{align*}
	\[
	Q_k = \left\{ \bigcup_{A \in S} A \mid S \text{ is a strongly connected component in } \Gamma_{k-1}(\mathcal A) \right\},
	\]
	and $E_{ \le k} = E_k \cup I$ where
	\begin{align*}
		E_k & =  \{ (C, C') \in Q_k \times Q_k \mid C \ne C',~  C \supseteq \excl{w}, 
		\\ & \qquad
		C' \cap \dupl{w} \ne \emptyset \text{ for some } w \in \Sigma^* \text{ with } \defect{w} = k \}, \\
		I   & =  \{ (C, C') \in Q_k \times Q_k \mid \exists (D, D') \in E_{k-1} : D \cap C \ne \emptyset \land 
		D' \cap C' \ne \emptyset \}.
	\end{align*}
	We call the edges in $E_k$ the edges \emph{induced}, or \emph{forced, by words of defect $k$} \IN{forced edges in $\Gamma_k(\mathcal A)$}
	and the edges in $I$ the edges \emph{inherited from $\Gamma_{k-1}$}. \IN{inherited edges in $\Gamma_k(\mathcal A)$}
\end{definition}

%In the following, mostly the digraph $\Gamma_1(\mathcal A) = (Q_1, E_1)$ will be used. 
Unraveling the definition,
the elements in $Q_1$ are singleton sets
and so can be identified with the elements in~$Q$. With respect to this identification,
which we will use in the following without special mention, we have $Q_1 = Q$ and 
\begin{equation}\label{syncmax:eqn:E1}
	E_1 = \{ (p, q) \mid \exists u \in \Sigma^* \text{ with $\excl{u} = \{p\}$ and $\dupl{u} = \{q\}$} \}.
\end{equation}

As an example for the construction, taken from~\cite{BondarCV23}, consider the automaton $\mathcal{E}_5$ with the state set $\{1,2,3,4,5\}$ and eight input letters, denoted $a_{[1]}$, $a_{[2]}$, $a_{[3]}$, $a_{[4]}$, $a_{[5]}$, $a_{[1,2]}$, $a_{[4,5]}$, and $a_{[1,3]}$
from \cref{fig:e5}. The transition function of $\mathcal{E}_5$ is given in Table~\ref{tb:e5}.
\begin{table}[ht]
	\caption{The transition table of the automaton $\mathcal{E}_5$}\label{tb:e5}
	\begin{center}
		\begin{tabular}{c@{\ }|@{\ }c@{\ }c@{\ }c@{\ }c@{\ }c@{\ }c@{\ }c@{\ }c}
			$q$  & $\delta(q, a_{[1]})$ & $\delta(q,a_{[2]})$ & $\delta(q, a_{[3]})$  & $\delta(q, a_{[4]})$ & $\delta(q, a_{[5]})$ & $\delta(q, a_{[1,2]})$ & $\delta(q, a_{[4,5]})$ & $\delta(q, a_{[1,3]})$ \mathstrut\\
			\hline
			1 & 2 & 1 & 1 & 1 & 1 & 3 & 1 & 4\mathstrut\\
			2 & 2 & 1 & 1 & 2 & 2 & 3 & 1 & 4\mathstrut\\
			3 & 3 & 3 & 2 & 3 & 3 & 3 & 2 & 4\mathstrut\\
			4 & 4 & 4 & 4 & 5 & 4 & 4 & 3 & 5\mathstrut\\
			5 & 5 & 5 & 5 & 5 & 4 & 5 & 3 & 5\mathstrut
		\end{tabular}
	\end{center}
\end{table}


The stages of the construction of $\Gamma_k(\mathcal E_5)$ are shown in \cref{sync_max:ex:compl_reach}. For example, to construct 
$\Gamma_1(\mathcal E_r)$, one has to compute the transformations caused by words of defect one with respect to $\mathcal{E}_5$.
A word of defect one in $\mathrsfs{E}_5$ must be a product of the defect-one letters $a_{[1]}$, $a_{[2]}$, $a_{[3]}$, $a_{[4]}$, $a_{[5]}$. The `multiplication table' for these words is given next, in which $*$ means that the defect of the product is larger than 1:
\begin{center}
	\begin{tabular}{ c| c | c | c | c | c |}
		& $a_{[1]}$ & $a_{[2]}$ & $a_{[3]}$ & $a_{[4]}$ & $a_{[5]}$ \\
		\hline
		$a_{[1]}$ & $a_{[1]}$ & $a_{[2]}$  & $a_{[3]}$ & $*$ & $*$\\
		\hline
		$a_{[2]}$ & $a_{[1]}$ & $a_{[2]}$  & $a_{[3]}$ & $*$ & $*$\\
		\hline
		$a_{[3]}$ & $*$ & $*$  & $*$ & $*$ & $*$\\
		\hline
		$a_{[4]}$ & $*$ & $*$  & $*$ & $a_{[4]}$ & $a_{[5]}$\\
		\hline
		$a_{[5]}$ & $*$ & $*$  & $*$ & $a_{[4]}$ & $a_{[5]}$\\
		\hline
	\end{tabular}
\end{center}

\begin{figure}[ht]
	\begin{center}
		\unitlength=0.95mm
		\begin{tikzpicture}[node distance = 1cm and 2.2cm]
			%\tikzstyle{vertex}=[circle,draw,minimum size=5pt,inner sep=1pt]
			\node[state] (1) {$1$};
			\node[state, left = of 1] (2) {$2$};
			\node[state, right = of 1] (3) {$3$};
			\node[state, right = of 3] (4) {$4$};
			\node[state, right = of 4] (5) {$5$};
			
			\path[->] 
			(1) edge [bend left=10,below] node {$a_{[1]}$} (2)
			(2) edge [above,bend left=10] node {$a_{[2]}, a_{[3]},a_{[4,5]}$} (1)
			(2) edge [bend right=35, above] node {$a_{[1,2]}$} (3)
			(3) edge [below,bend left=55] node {$a_{[3]}$,$a_{[4,5]}$} (2);	 
			
			\path[->]
			(2) edge [bend left=60,above] node {$a_{[1,3]}$} (4)
			(1) edge [bend left=40,above] node {$a_{[1,3]}$} (4)
			(4) edge [bend left=10,above] node {$a_{[4]}, a_{[1,3]}$} (5)
			(5) edge [below, bend left=10] node {$a_{[5]}$} (4)
			(3) edge [bend left=10, above] node {$a_{[1,3]}$} (4)
			(4) edge [bend left=10, below] node {$a_{[4,5]}$} (3)
			(5) edge [bend left, below] node {$a_{[4,5]}$} (3)
			(1) edge [left,above] node {$a_{[1,2]}$} (3);
		\end{tikzpicture}
		\caption[Automaton $\mathcal E_5$ to illustrate graph constructions $\Gamma_k(\mathcal E_5)$]{The automaton $\mathcal{E}_5$ from~\cite{BondarCV23} with self-loops omitted. Here $\excl{a_{[i]}} = \{i\}$ and $\excl{a_{[i,j]}} = \{i,i+1,\ldots,j\}$ for $1 \le i < j \le 5$.} \label{fig:e5}
	\end{center}
\end{figure}
\begin{figure}[ht]
	\begin{minipage}{0.48\textwidth}
		\begin{subfigure}{1.0\textwidth}\centering
			\begin{tikzpicture}[node distance = 1cm and 1.8cm]
				\tikzstyle{vertex}=[]
				\node[vertex] (1) {$\{1\}$};
				\node[vertex, left = of 1] (2)  {$\{2\}$};
				\node[vertex, right = of 1] (3) {$\{3\}$};
				\node[vertex, below right = of 2, xshift=-1.35cm] (4) {$\{4\}$};
				\node[vertex, right = of 4] (5) {$\{5\}$};
				
				\path[->, dashed] (1) edge [below] node {$a_{[1]}$} (2)
				(2) edge [bend left, above] node {$a_{[2]}$} (1)
				(3) edge [above] node {$a_{[3]}$} (1)
				(4) edge [bend left, above] node {$a_{[4]}$} (5)
				(5) edge [below] node {$a_{[5]}$} (4);
			\end{tikzpicture}
			\caption{The graph $\Gamma_1(\mathcal A)$.}
		\end{subfigure}	
	\end{minipage}\hfill
	\begin{minipage}{0.48\textwidth}
		\begin{subfigure}{1.0\textwidth}\centering
			\begin{tikzpicture}[node distance = 1cm and 1.8cm]
				\tikzstyle{vertex}=[]
				\node[vertex] (1) {$\{3\}$};
				\node[vertex, left = of 1] (2) {$\{1,2\}$};
				\node[vertex, right = of 1] (3) {$\{4,5\}$};
				
				\path[->] (1) edge (2)
				(2) edge [dashed,bend left,above] node {$a_{[1,2]}$} (1)
				(3) edge [dashed,bend left,above] node {$a_{[4,5]}$} (2)
				(3) edge [dashed,above] node {$a_{[4,5]}$} (1);
			\end{tikzpicture}
			\caption{The graph $\Gamma_2(\mathcal A)$.}
		\end{subfigure}
		\begin{subfigure}{1.0\textwidth}\centering
			\medskip 
			\begin{tikzpicture}[node distance = 1cm and 1.8cm]
				\tikzstyle{vertex}=[]
				\node[vertex] (1) {$\{4,5\}$};
				\node[vertex, left = of 1] (2) {$\{1,2,3\}$};
				
				\path[->] (2) edge [dashed, bend left, above] node {$a_{[1,3]}$} (1)
				(1) edge (2);
			\end{tikzpicture}
			\caption{The graph $\Gamma_3(\mathcal A)$.}
		\end{subfigure}
	\end{minipage} 
	\caption[Example for graphs $\Gamma_k$.]{The graphs $\Gamma_k(\mathcal A)$ for $k \in \{1,2,3\}$. Dashed edges are induced by words
		of defect $1$, $2$, or $3$ written above the edge. Solid edges are inherited
		from the previous graph.}
	\label{sync_max:ex:compl_reach}
\end{figure}


The graph $\Gamma_1(\mathcal A)$ was introduced in~\cite{BondarV16}
to give sufficient condition for complete reachability.
Later, all the graphs $\Gamma_k(\mathcal A)$ were introduced in~\cite{BondarV18,BondarCV23} to characterize 
complete reachability as stated in the following theorem.

\begin{theorem}[Bondar et al.~\cite{BondarV18,BondarCV23}]
	\label{sync_max:thm:char_compl_reach}
%	If there exists $0 < k < n$ such $\Gamma_k(\mathcal A)$
%	is strongly connected, then $\mathcal A$ is completely reachable
%	and every non-empty subset is reachable by a product of words of
%	defect at most $k$.
%	Otherwise, there exists $0 < k < n$
%	such that each subset in $Q_k$ contains less than $k$
%	elements and some subset of $Q$ with at least $n - k$
%	states is not reachable in $\mathcal A$.	
	The automaton $\mathcal A$ is completely reachable 
	if and only if there exists 
	$0 < k < n$ such that $\Gamma_k(\mathcal A)$ is strongly connected (and in this case every non-empty subset
	is reachable by a product of words of defect at most~$k$).
\end{theorem}

%\begin{figure}[ht]
%	\begin{minipage}{0.6\textwidth}
%			\begin{tikzpicture}[node distance = 1.5cm and 2.7cm]
%			\tikzstyle{vertex}=[circle,draw,minimum size=5pt,inner sep=1pt]
%		\node[vertex] (1) {$1$};
%		\node[vertex, left = of 1] (2) {$2$};
%		\node[vertex, below = of 1, yshift=-1.0cm] (3) {$3$};
%		\node[vertex, above right = of 1] (4) {$4$};
%		\node[vertex, below = of 4, yshift=-3cm] (5) {$5$};
%		
%		\path[->] 
%		(1) edge [below] node {$a_{[1]}$} (2)
%		(2) edge [above,bend left] node {$a_{[2]}, a_{[3]},a_{[4,5]}$} (1)
%		(2) edge [left,pos=.5] node {$a_{[2,3]}$} (3)
%		(3) edge [below,bend left, out=50,in=120,pos=.5] node [align=left] {\parbox{1cm}{$a_{[3]}$ \newline $a_{[4,5]}$}} (2);	 
%		
%		\path[->]
%		(2) edge [bend left,out=90,above] node {$a_{[1,3]}$} (4)
%		(1) edge [bend left,above] node {$a_{[1,3]}$} (4)
%		(4) edge [bend left,right,pos=.2,out=50,in=130] node {$a_{[4]}, a_{[1,3]}$} (5)
%		(5) edge [right,pos=.5] node {$a_{[5]}$} (4)
%		(3) edge [left,pos=0.8] node {$a_{[1,3]}$} (4)
%		(4) edge [bend left,pos=0.7,right] node {$a_{[4,5]}$} (3)
%		(5) edge [below] node {$a_{[4,5]}$} (3)
%		(1) edge [left, bend left] node {$a_{[1,2]}$} (3);
%		
%				\end{tikzpicture}
%		\end{minipage}\hfill
%	\begin{minipage}{0.35\textwidth}
%			\begin{subfigure}{1.0\textwidth}
%					\begin{tikzpicture}[node distance = 1cm and 0.5cm]
%							\tikzstyle{vertex}=[]
%							\node[vertex] (1) {$\{1\}$};
%							\node[vertex, left = of 1] (2)  {$\{2\}$};
%							\node[vertex, right = of 1] (3) {$\{3\}$};
%							\node[vertex, right = of 3] (4) {$\{4\}$};
%							\node[vertex, below = of 4] (5) {$\{5\}$};
%							
%							\path[->, dashed] (1) edge (2)
%							(2) edge [bend left] (1)
%							(3) edge (1)
%							(4) edge [bend right] (5)
%							(5) edge (4);
%						\end{tikzpicture}
%					\caption{The graph $\Gamma_1(\mathcal A)$.}
%				\end{subfigure}
%			\begin{subfigure}{1.0\textwidth}
%					\medskip 	
%					\begin{tikzpicture}[node distance = 1cm and 0.5cm]
%							\tikzstyle{vertex}=[]
%							\node[vertex] (1) {$\{3\}$};
%							\node[vertex, left = of 1] (2) {$\{1,2\}$};
%							\node[vertex, right = of 1] (3) {$\{4,5\}$};
%							
%							\path[->] (1) edge (2)
%							(2) edge [dashed,bend left] (1)
%							(3) edge [dashed,bend left] (2)
%							(3) edge [dashed] (1);
%						\end{tikzpicture}
%					\caption{The graph $\Gamma_2(\mathcal A)$.}
%				\end{subfigure}
%			\begin{subfigure}{1.0\textwidth}
%					\medskip 
%					\begin{tikzpicture}[node distance = 1cm and 0.5cm]
%							\tikzstyle{vertex}=[]
%							\node[vertex] (1) {$\{\{4,5\}\}$};
%							\node[vertex, left = of 1] (2) {$\{\{1,2\},\{3\}\}\}$};
%							
%							\path[->] (2) edge [dashed, bend left] (1)
%							(1) edge (2);
%						\end{tikzpicture}
%					\caption{The graph $\Gamma_3(\mathcal A)$.}
%				\end{subfigure}
%			
%		\end{minipage} 
%	\caption[Example for graphs $\Gamma_k$.]{Left: $\semiautomat{ E_5}{\{1,2,3,4,5\}}{\{a_{[1]}, a_{[2]}, a_{[3]}, a_{[4]}, a_{[5]}, a_{[1,2]}, a_{[1,3]}, a_{[4,5]}\}}{\delta}$
%			taken from~\cite{DBLP:journals/corr/abs-2201-05075} (self-loops omitted). Right: The graphs $\Gamma_k(\mathcal A)$ for $k \in \{1,2,3\}$. Dashed edges are induced by words
%			of defect $1$, $2$ and $3$ respectively, the solid edges are inherited
%			from the previous graph.}
%	\label{sync_max:ex:compl_reach}
%\end{figure}

The following proposition states that the vertices of each $\Gamma_k(\mathcal A)$
partition~$Q$ and the partition given by $Q_{k-1}$ is coarser than the partition
given by $Q_k$ for~$k \ge 1$. 

\begin{proposition} \label{sync_max:prop:BVCgraphs-vertices}
Let $\semiautomat{A}{Q}{\Sigma}{\delta}$ and $\Gamma_k(\mathcal A) = (Q_k, E_k)$ for $k \ge 0$.
%be the Bondar-Volkov-Casas graphs.
Suppose $D, D' \in Q_{k}$ and $C \in Q_{k+1}$.
Then the following is true.
	\begin{enumerate}
		\item \label{sync_max:prop:BVCgraphs-vertices:1} $D = D'$ or $D \cap D' = \emptyset$.		
		
		\item \label{sync_max:prop:BVCgraphs-vertices:2} $D \cap C \ne \emptyset$ if and only if $D \subseteq C$.
	\end{enumerate}
\end{proposition}
\begin{proof}
\ref{sync_max:prop:BVCgraphs-vertices:1}: For $k = 0$ this is obvious. So, assume $k \ge 1$.
	Suppose there exists $p \in D \cap D'$ % B statt D, da davor!
	and let $q \in D$.
	By definition, we have $p \in B$, $p \in B'$
	and $q \in B''$ for $B, B', B'' \in Q_{k-1}$ with 
	$B \subseteq D$, $B' \subseteq D'$, $B'' \subseteq D$
	and where $B$ and $B''$ are in the same strongly connected component in $\Gamma_{k-1}(\mathcal A)$.
	Inductively, as $p \in B \cap B'$,
	we find $B = B'$.
	So, $B'$ and $B''$ are in the same strongly connected component in $\Gamma_{k-1}(\mathcal A)$
	and as $B' \subseteq D'$ we get $B'' \subseteq D'$
	so that $q \in D'$.
	Similarly, we find $D' \subseteq D$.	
	
\ref{sync_max:prop:BVCgraphs-vertices:2}: Suppose $q \in D \cap C$ for some $q \in Q$.
	As $q \in C$ there exists $D'' \in Q_{k-1}$
	with $q \in D''$ and $D'' \subseteq C$.
	By Item~\ref{sync_max:prop:BVCgraphs-vertices:1} we get $D = D''$, which gives the claim.	
\end{proof}

By \cref{sync_max:prop:BVCgraphs-vertices:2} of \cref{sync_max:prop:BVCgraphs-vertices}
the set $I$ in Definition~\ref{sync_max:def:BVCgraphs} 
can also be defined by the conditions $D \subseteq C$ and $D' \subseteq C'$.

Lastly, we introduce a variant of $\Gamma_1(\mathcal A)$ that we need later to establish bounds on shortest
words reaching non-empty subsets in classes of automata.

\label{def:Gamma1i-graph}
For $i \ge 1$, let $\Gamma_{1,i}(\mathcal A)$ be the subgraph of $\Gamma_1(\mathcal A)$
whose edges are all induced be words of defect one of the form $au$
with $a \in \Sigma$ being a letter of defect one and $u$ being a permutational word
with $|u| < i$, i.e., its edge set has the form
\begin{multline*}
\{ (p,q) \mid \excl{au} = \{p\} \text{ and }  \dupl{au} = \{q\} \\ 
\text{ for some $a \in \Sigma$ of defect one and a permutational word $u \in \Sigma^{< i}$} \}
\end{multline*}

\begin{theorem} \label{syncmax:thm:Gamma1-SCC-subgraph}
	Let $\semiautomat{A}{Q}{\Sigma}{\delta}$ be an $n$-state automaton.
	Suppose $\Gamma_{1,i}(\mathcal A)$ is strongly connected for $i \ge 1$.
	Then~$\mathcal A$ is completely reachable,
	every non-empty subset of size~$r$ is reachable by a word
	of length at most $i \cdot (n - r)$
	and a shortest synchronizing word has length at most $1 + i \cdot (n - 2)$.
\end{theorem}
\begin{proof}
	We can assume $n \ge 2$.
	Let $S \subseteq Q$ be a non-empty $r$-subset. First, we prove that~$S$ is reachable
	by a word of length at most $i \cdot ( n - r )$
	by induction on $k = |Q \setminus S|$.
	If $k = 0$, then $S = Q$
	and so~$S$ is reachable via the empty word
	and the claim holds.	
	Now let $k > 0$ so that~$S$ is a proper subset of~$Q$.
	Since $\Gamma_{1,i}(\mathcal A)$ is strongly connected,
	by \cref{prelim:lem:partition_SCC}
	it has an edge $(p, q)$ with $p \in Q \setminus S$ and $q \in S$.
	By the definition of $\Gamma_{1,i}(\mathcal A)$,
	there exists a word of defect one of the form $au$
	for some letter of defect one $a \in \Sigma$
	and a permutational word $u$ with $|u| < i$
	such that $\excl{au} = \{ p \}$ and $\dupl{au} = \{ q \}$.
	Since the excluded $p$ of~$au$ does not belong to~$S$
	we have $|\delta^{-1}(s, au)| \ge 1$
	for each $s \in S$
	where $|\delta^{-1}(q, au)| = 2$.
	Thus for $T = \delta^{-1}(S, w)$
	we have $|T| = |S| + 1$ and $\delta(T, au) = S$.
	Inductively, 
	$T$ is reachable by a word $w$ 
	of length at most $i \cdot ( n - |T| )$.
	Then $\delta(Q, wau) = S$
	and $|wau| \le i \cdot ( n - |T| ) + i = i \cdot ( n - |S| )$.
	This gives the first part of the claim by induction.
	
	Lastly, since~$\mathcal A$ is synchronizing,
	there exist distinct states $p, q \in Q$
	and a letter $a \in \Sigma$ with $\delta(p, a) = \delta(q, a)$.
	Now the set $\{p,q\}$ is reachable by a word $w$ of length at most $i \cdot (n - 2)$.
	Then the word $wa$ synchronizes~$\mathcal A$
	and has length at most $1 + i \cdot ( n - 2 )$.
\end{proof}